\documentclass[11pt]{article}
\newcommand{\ListaTitulo}{Lista 02 --- Modelos Lineares (Parte 1)}
% Preâmbulo compartilhado — MATD48, Listas de Exercícios 2026
% Usado por Lista{NN}.tex e Gabarito{NN}.tex via % Preâmbulo compartilhado — MATD48, Listas de Exercícios 2026
% Usado por Lista{NN}.tex e Gabarito{NN}.tex via % Preâmbulo compartilhado — MATD48, Listas de Exercícios 2026
% Usado por Lista{NN}.tex e Gabarito{NN}.tex via \input{preamble}

\usepackage[utf8]{inputenc}
\usepackage[T1]{fontenc}
\usepackage[brazil]{babel}
\usepackage{fullpage}
\usepackage{amsmath,amssymb}
\usepackage{enumitem}
\usepackage{xcolor}
\usepackage{fancyhdr}
\usepackage{titlesec}
\usepackage{listings}
\usepackage{hyperref}
\usepackage{booktabs}
\usepackage{graphicx}

\definecolor{matd48azul}{HTML}{0B4F6C}
\definecolor{matd48verde}{HTML}{2E8B57}

\titleformat{\section}{\normalfont\Large\bfseries\color{matd48azul}}{\thesection}{1em}{}
\titleformat{\subsection}{\normalfont\large\bfseries\color{matd48azul}}{\thesubsection}{1em}{}

\providecommand{\ListaTitulo}{Lista de Exercícios}

\setlength{\headheight}{14pt}
\pagestyle{fancy}
\fancyhf{}
\lhead{\small MATD48 --- Planejamento de Experimentos A}
\rhead{\small \ListaTitulo}
\cfoot{\thepage}
\renewcommand{\headrulewidth}{0.4pt}

\lstset{
  basicstyle=\ttfamily\small,
  breaklines=true,
  frame=single,
  columns=fullflexible,
  backgroundcolor=\color{gray!8},
  commentstyle=\color{matd48verde},
  keywordstyle=\color{matd48azul}\bfseries,
  language=R,
  extendedchars=true,
  literate=%
    {á}{{\'a}}1 {à}{{\`a}}1 {â}{{\^a}}1 {ã}{{\~a}}1
    {é}{{\'e}}1 {ê}{{\^e}}1 {í}{{\'i}}1 {ó}{{\'o}}1
    {ô}{{\^o}}1 {õ}{{\~o}}1 {ú}{{\'u}}1 {ç}{{\c c}}1
    {Á}{{\'A}}1 {À}{{\`A}}1 {Â}{{\^A}}1 {Ã}{{\~A}}1
    {É}{{\'E}}1 {Ê}{{\^E}}1 {Í}{{\'I}}1 {Ó}{{\'O}}1
    {Ô}{{\^O}}1 {Õ}{{\~O}}1 {Ú}{{\'U}}1 {Ç}{{\c C}}1
}

\hypersetup{colorlinks=true, linkcolor=matd48azul, urlcolor=matd48azul, citecolor=matd48azul}

\newcommand{\cabecalho}[2]{%
  \begin{center}
    {\Large\bfseries\color{matd48azul} #1}\\[0.3em]
    {\large #2}\\[0.3em]
    \rule{\linewidth}{0.6pt}
  \end{center}
  \vspace{0.5em}
}

\newenvironment{questao}[1]{\vspace{1em}\noindent\textbf{Questão #1.}\hspace{0.4em}}{\par}
\newenvironment{solucao}{\vspace{0.3em}\noindent\textit{Solução.}\hspace{0.4em}\color{black!85}}{\par\vspace{0.5em}}


\usepackage[utf8]{inputenc}
\usepackage[T1]{fontenc}
\usepackage[brazil]{babel}
\usepackage{fullpage}
\usepackage{amsmath,amssymb}
\usepackage{enumitem}
\usepackage{xcolor}
\usepackage{fancyhdr}
\usepackage{titlesec}
\usepackage{listings}
\usepackage{hyperref}
\usepackage{booktabs}
\usepackage{graphicx}

\definecolor{matd48azul}{HTML}{0B4F6C}
\definecolor{matd48verde}{HTML}{2E8B57}

\titleformat{\section}{\normalfont\Large\bfseries\color{matd48azul}}{\thesection}{1em}{}
\titleformat{\subsection}{\normalfont\large\bfseries\color{matd48azul}}{\thesubsection}{1em}{}

\providecommand{\ListaTitulo}{Lista de Exercícios}

\setlength{\headheight}{14pt}
\pagestyle{fancy}
\fancyhf{}
\lhead{\small MATD48 --- Planejamento de Experimentos A}
\rhead{\small \ListaTitulo}
\cfoot{\thepage}
\renewcommand{\headrulewidth}{0.4pt}

\lstset{
  basicstyle=\ttfamily\small,
  breaklines=true,
  frame=single,
  columns=fullflexible,
  backgroundcolor=\color{gray!8},
  commentstyle=\color{matd48verde},
  keywordstyle=\color{matd48azul}\bfseries,
  language=R,
  extendedchars=true,
  literate=%
    {á}{{\'a}}1 {à}{{\`a}}1 {â}{{\^a}}1 {ã}{{\~a}}1
    {é}{{\'e}}1 {ê}{{\^e}}1 {í}{{\'i}}1 {ó}{{\'o}}1
    {ô}{{\^o}}1 {õ}{{\~o}}1 {ú}{{\'u}}1 {ç}{{\c c}}1
    {Á}{{\'A}}1 {À}{{\`A}}1 {Â}{{\^A}}1 {Ã}{{\~A}}1
    {É}{{\'E}}1 {Ê}{{\^E}}1 {Í}{{\'I}}1 {Ó}{{\'O}}1
    {Ô}{{\^O}}1 {Õ}{{\~O}}1 {Ú}{{\'U}}1 {Ç}{{\c C}}1
}

\hypersetup{colorlinks=true, linkcolor=matd48azul, urlcolor=matd48azul, citecolor=matd48azul}

\newcommand{\cabecalho}[2]{%
  \begin{center}
    {\Large\bfseries\color{matd48azul} #1}\\[0.3em]
    {\large #2}\\[0.3em]
    \rule{\linewidth}{0.6pt}
  \end{center}
  \vspace{0.5em}
}

\newenvironment{questao}[1]{\vspace{1em}\noindent\textbf{Questão #1.}\hspace{0.4em}}{\par}
\newenvironment{solucao}{\vspace{0.3em}\noindent\textit{Solução.}\hspace{0.4em}\color{black!85}}{\par\vspace{0.5em}}


\usepackage[utf8]{inputenc}
\usepackage[T1]{fontenc}
\usepackage[brazil]{babel}
\usepackage{fullpage}
\usepackage{amsmath,amssymb}
\usepackage{enumitem}
\usepackage{xcolor}
\usepackage{fancyhdr}
\usepackage{titlesec}
\usepackage{listings}
\usepackage{hyperref}
\usepackage{booktabs}
\usepackage{graphicx}

\definecolor{matd48azul}{HTML}{0B4F6C}
\definecolor{matd48verde}{HTML}{2E8B57}

\titleformat{\section}{\normalfont\Large\bfseries\color{matd48azul}}{\thesection}{1em}{}
\titleformat{\subsection}{\normalfont\large\bfseries\color{matd48azul}}{\thesubsection}{1em}{}

\providecommand{\ListaTitulo}{Lista de Exercícios}

\setlength{\headheight}{14pt}
\pagestyle{fancy}
\fancyhf{}
\lhead{\small MATD48 --- Planejamento de Experimentos A}
\rhead{\small \ListaTitulo}
\cfoot{\thepage}
\renewcommand{\headrulewidth}{0.4pt}

\lstset{
  basicstyle=\ttfamily\small,
  breaklines=true,
  frame=single,
  columns=fullflexible,
  backgroundcolor=\color{gray!8},
  commentstyle=\color{matd48verde},
  keywordstyle=\color{matd48azul}\bfseries,
  language=R,
  extendedchars=true,
  literate=%
    {á}{{\'a}}1 {à}{{\`a}}1 {â}{{\^a}}1 {ã}{{\~a}}1
    {é}{{\'e}}1 {ê}{{\^e}}1 {í}{{\'i}}1 {ó}{{\'o}}1
    {ô}{{\^o}}1 {õ}{{\~o}}1 {ú}{{\'u}}1 {ç}{{\c c}}1
    {Á}{{\'A}}1 {À}{{\`A}}1 {Â}{{\^A}}1 {Ã}{{\~A}}1
    {É}{{\'E}}1 {Ê}{{\^E}}1 {Í}{{\'I}}1 {Ó}{{\'O}}1
    {Ô}{{\^O}}1 {Õ}{{\~O}}1 {Ú}{{\'U}}1 {Ç}{{\c C}}1
}

\hypersetup{colorlinks=true, linkcolor=matd48azul, urlcolor=matd48azul, citecolor=matd48azul}

\newcommand{\cabecalho}[2]{%
  \begin{center}
    {\Large\bfseries\color{matd48azul} #1}\\[0.3em]
    {\large #2}\\[0.3em]
    \rule{\linewidth}{0.6pt}
  \end{center}
  \vspace{0.5em}
}

\newenvironment{questao}[1]{\vspace{1em}\noindent\textbf{Questão #1.}\hspace{0.4em}}{\par}
\newenvironment{solucao}{\vspace{0.3em}\noindent\textit{Solução.}\hspace{0.4em}\color{black!85}}{\par\vspace{0.5em}}


\begin{document}

\cabecalho{Lista de Exercícios 02}{Modelo linear, mínimos quadrados e estimabilidade (Capítulo 2 / Aula 02)}

\noindent \textit{Entrega: até a próxima aula. Justifique todas as respostas — nestas listas, a
justificativa vale mais que a resposta final. O gabarito completo está em \texttt{Gabarito02.pdf}.}

\begin{questao}{1} \textbf{(Ciência de dados --- mínimos quadrados na mão)}

Uma pequena amostra de 5 medições relaciona o número de \emph{scripts} carregados por uma página
web ($x$) ao tempo de carregamento em décimos de segundo ($y$):

\begin{center}
\begin{tabular}{c|ccccc}
$x$ & 2 & 4 & 6 & 8 & 10 \\
\hline
$y$ & 15 & 22 & 25 & 33 & 38
\end{tabular}
\end{center}

Considere o modelo $y_i = \beta_0 + \beta_1 x_i + \varepsilon_i$.

\begin{enumerate}[label=(\alph*)]
  \item Escreva explicitamente a matriz de delineamento $\mathbf{X}$ ($5 \times 2$) e o vetor
    $\mathbf{Y}$ para este conjunto de dados.
  \item Calcule $\mathbf{X}'\mathbf{X}$ e $\mathbf{X}'\mathbf{Y}$.
  \item Resolva as equações normais $\mathbf{X}'\mathbf{X}\hat{\boldsymbol\beta} =
    \mathbf{X}'\mathbf{Y}$ para obter $\hat{\beta}_0$ e $\hat{\beta}_1$. (Dica:
    $(\mathbf{X}'\mathbf{X})^{-1}$ para uma matriz $2\times 2$
    $\begin{bmatrix} a & b \\ b & c\end{bmatrix}$ é $\frac{1}{ac-b^2}\begin{bmatrix} c & -b \\ -b &
    a \end{bmatrix}$.)
  \item Interprete $\hat{\beta}_1$ no contexto do problema.
  \item Sem recalcular tudo: qual deveria ser a soma dos resíduos $\sum_i (y_i - \hat y_i)$?
    Justifique sua resposta a partir das equações normais, sem apenas calcular os resíduos.
\end{enumerate}
\end{questao}

\begin{questao}{2} \textbf{(Psicologia --- estimabilidade)}

Um estudo mede o tempo de reação (ms) de 6 participantes, 2 em cada um de três grupos
experimentais ($G_1$, $G_2$, $G_3$). O modelo utilizado, na parametrização
\emph{sobreparametrizada} usual de um fator, é

$$
y_{ij} = \mu + \alpha_j + \varepsilon_{ij}, \qquad j = 1,2,3, \quad i = 1,2,
$$

com $\mathbf{X}$ formada por uma coluna de 1's (para $\mu$) e uma coluna indicadora para cada
grupo (para $\alpha_1$, $\alpha_2$, $\alpha_3$).

\begin{enumerate}[label=(\alph*)]
  \item Escreva a matriz $\mathbf{X}$ ($6 \times 4$) explicitamente, com colunas na ordem $(\mu,
    \alpha_1, \alpha_2, \alpha_3)$.
  \item Qual é o posto de $\mathbf{X}$? Explique por que ele é menor do que o número de colunas.
  \item Para cada função a seguir, diga se é \textbf{estimável} e justifique usando a definição
    (existe $\mathbf{a}$ tal que $\mathbf{a}'\mathbf{X} = \boldsymbol\lambda'$?):
    \begin{enumerate}[label=(\roman*)]
      \item $\mu$
      \item $\mu + \alpha_1$
      \item $\alpha_1 - \alpha_2$
      \item $\alpha_1 + \alpha_2 + \alpha_3$
    \end{enumerate}
\end{enumerate}
\end{questao}

\begin{questao}{3} \textbf{(Agricultura --- especificação do modelo)}

Uma equipe agronômica testa o efeito de dois fatores sobre a produtividade do milho: dose de
fertilizante (baixa, alta) e variedade (V1, V2), com uma repetição de cada combinação (4 parcelas
no total). Usando a codificação padrão do \texttt{R} (contrastes de tratamento, que descartam o
nível de referência de cada fator), o modelo é

$$
y_i = \beta_0 + \beta_1 \, \mathbf{1}\{\text{dose alta}\} + \beta_2\, \mathbf{1}\{\text{V2}\} +
\varepsilon_i.
$$

\begin{enumerate}[label=(\alph*)]
  \item Escreva $\mathbf{X}$ ($4 \times 3$) para as quatro parcelas (uma de cada combinação
    dose$\times$variedade).
  \item $\mathbf{X}$ tem posto coluna completo? Justifique calculando (ou argumentando sobre) o
    posto.
  \item Interprete $\beta_0$, $\beta_1$ e $\beta_2$ em termos das médias esperadas das quatro
    combinações de tratamento.
\end{enumerate}
\end{questao}

\begin{questao}{4} \textbf{(Ciência de dados --- geometria dos mínimos quadrados)}

Mostre, a partir das equações normais $\mathbf{X}'\mathbf{X}\hat{\boldsymbol\beta} =
\mathbf{X}'\mathbf{Y}$, que:

\begin{enumerate}[label=(\alph*)]
  \item $\mathbf{X}'\mathbf{e} = \mathbf{0}$, em que $\mathbf{e} = \mathbf{Y} -
    \mathbf{X}\hat{\boldsymbol\beta}$ é o vetor de resíduos. (Dica: reescreva as equações normais
    isolando $\mathbf{X}'(\mathbf{Y} - \mathbf{X}\hat{\boldsymbol\beta})$.)
  \item Se a primeira coluna de $\mathbf{X}$ é um vetor de 1's (intercepto), conclua a partir do
    item (a) que $\sum_i e_i = 0$.
  \item Isso significa que o modelo "acerta em média" no sentido de não haver nenhum padrão
    sistemático não capturado nos resíduos? Justifique com um contraexemplo conceitual (não
    precisa ser numérico).
\end{enumerate}
\end{questao}

\begin{questao}{5} \textbf{(Ciência de dados --- exercício computacional em R)}

Considere novamente os dados da Questão 1. O código abaixo está incompleto.

\begin{lstlisting}
x <- c(2, 4, 6, 8, 10)
y <- c(15, 22, 25, 33, 38)
X <- cbind(1, x)

# (a) complete a linha abaixo para obter beta_hat via algebra de matrizes
beta_hat <- ______________________

# (b) obtenha os valores ajustados e os residuos
y_hat <- X %*% beta_hat
residuos <- ______________________

# (c) confirme a Questao 1(e): a soma dos residuos deve ser (numericamente) zero
sum(residuos)
\end{lstlisting}

\begin{enumerate}[label=(\alph*)]
  \item Complete as duas linhas indicadas.
  \item Rode o código (mentalmente ou em sua máquina) e confirme que os valores de
    $\hat{\boldsymbol\beta}$ batem com os obtidos manualmente na Questão 1.
  \item Modifique o código para comparar sua estimativa com a saída de \texttt{lm(y \textasciitilde{}
    x)}. As duas devem coincidir --- explique, em uma frase, por quê.
\end{enumerate}
\end{questao}

\begin{questao}{6} \textbf{(Dedução --- equações normais)}

Considere o modelo linear geral $\mathbf{Y} = \mathbf{X}\boldsymbol\beta + \boldsymbol\varepsilon$
e a soma de quadrados dos resíduos $S(\boldsymbol\beta) = (\mathbf{Y} -
\mathbf{X}\boldsymbol\beta)'(\mathbf{Y} - \mathbf{X}\boldsymbol\beta)$.

\begin{enumerate}[label=(\alph*)]
  \item Expanda $S(\boldsymbol\beta)$ em termos de $\mathbf{Y}'\mathbf{Y}$,
    $\boldsymbol\beta'\mathbf{X}'\mathbf{Y}$ e $\boldsymbol\beta'\mathbf{X}'\mathbf{X}\boldsymbol\beta$.
  \item Usando as regras de derivação matricial $\partial(\mathbf{a}'\boldsymbol\beta)/\partial
    \boldsymbol\beta = \mathbf{a}$ e $\partial(\boldsymbol\beta'\mathbf{A}\boldsymbol\beta)/\partial
    \boldsymbol\beta = 2\mathbf{A}\boldsymbol\beta$ (para $\mathbf{A}$ simétrica), calcule
    $\partial S(\boldsymbol\beta)/\partial \boldsymbol\beta$.
  \item Iguale a derivada a $\mathbf{0}$ e obtenha as equações normais
    $\mathbf{X}'\mathbf{X}\hat{\boldsymbol\beta} = \mathbf{X}'\mathbf{Y}$.
  \item Explique por que a matriz Hessiana ($2\mathbf{X}'\mathbf{X}$) sendo semidefinida positiva
    garante que a solução das equações normais é um \emph{mínimo} de $S(\boldsymbol\beta)$, e não
    um máximo ou ponto de sela.
\end{enumerate}
\end{questao}

\begin{questao}{7} \textbf{(Dedução/computacional --- pseudoinversa de Moore-Penrose via SVD)}

Um experimento com dois tratamentos, cada um com 2 réplicas, produz o modelo sobreparametrizado
$y_{ij} = \mu + \tau_i + \varepsilon_{ij}$ ($i=1,2$), com

$$
\mathbf{X} = \begin{bmatrix} 1&1&0\\1&1&0\\1&0&1\\1&0&1 \end{bmatrix}, \qquad
\mathbf{y} = \begin{bmatrix} 15{,}2\\14{,}3\\13{,}5\\14{,}5 \end{bmatrix}, \qquad
\mathbf{A} := \mathbf{X}'\mathbf{X} = \begin{bmatrix} 4&2&2\\2&2&0\\2&0&2 \end{bmatrix}.
$$

Como $\mathbf{A}$ é simétrica, sua decomposição em valores singulares coincide com sua
decomposição espectral, $\mathbf{A} = \mathbf{V}\mathbf{D}\mathbf{V}'$, com autovalores/autovetores

$$
\lambda_1 = 6,\ \mathbf{v}_1 = \tfrac{1}{\sqrt6}(2,1,1)'; \qquad
\lambda_2 = 2,\ \mathbf{v}_2 = \tfrac{1}{\sqrt2}(0,-1,1)'; \qquad
\lambda_3 = 0,\ \mathbf{v}_3 = \tfrac{1}{\sqrt3}(1,-1,-1)'.
$$

\begin{enumerate}[label=(\alph*)]
  \item Verifique, multiplicando $\mathbf{A}\mathbf{v}_3$, que $\mathbf{v}_3$ é de fato um
    autovetor de autovalor $0$ (isto é, $\mathbf{A}\mathbf{v}_3 = \mathbf{0}$) --- e explique por
    que a existência de um autovalor nulo já era esperada, dado o posto incompleto de
    $\mathbf{X}$.
  \item Escreva a matriz $\mathbf{D}^+$ (pseudoinversa de Moore-Penrose de $\mathbf{D} =
    \mathrm{diag}(6,2,0)$): inverta apenas os valores singulares não nulos e mantenha $0$ na
    posição correspondente a $\lambda_3$.
  \item A pseudoinversa é $\mathbf{A}^+ = \mathbf{V}\mathbf{D}^+\mathbf{V}'$, com $\mathbf{V} =
    [\mathbf{v}_1\ \mathbf{v}_2\ \mathbf{v}_3]$. Calcule explicitamente a entrada $(1,1)$ de
    $\mathbf{A}^+$ (isto é, $\mathbf{e}_1'\mathbf{A}^+\mathbf{e}_1$), mostrando as contas com
    $\mathbf{v}_1, \mathbf{v}_2$ (o termo em $\mathbf{v}_3$ some, porque é multiplicado por $0$ em
    $\mathbf{D}^+$). Confira que sua resposta bate com $1/9 \approx 0{,}1111$.
  \item Complete o código abaixo, que calcula $\mathbf{A}^+$ inteira via \texttt{svd()} e a compara
    com \texttt{MASS::ginv(A)}:
\begin{lstlisting}
A <- matrix(c(4,2,2, 2,2,0, 2,0,2), ncol = 3, byrow = TRUE)
sv <- svd(A)
D_mais <- diag(______________________)      # (i) inverta so os valores singulares > 1e-8
A_mais <- ______________________             # (ii) V %*% D+ %*% t(U)
max(abs(A_mais - MASS::ginv(A)))
\end{lstlisting}
  \item Use $\mathbf{A}^+$ para obter $\mathbf{b} = \mathbf{A}^+\mathbf{X}'\mathbf{y}$ e, a partir
    de $\mathbf{b}$, as funções estimáveis $\mu+\tau_1$, $\mu+\tau_2$ e $\tau_1-\tau_2$
    (multiplicando $\mathbf{b}$ pelos vetores $\boldsymbol\lambda'$ apropriados, como na Questão
    2). Se algum colega usar uma inversa generalizada \emph{diferente} de $\mathbf{A}^+$ para
    resolver as equações normais, seu $\mathbf{b}$ será diferente do seu --- mas os três valores
    estimáveis calculados devem coincidir. Explique, em uma frase, por que a teoria garante essa
    coincidência mesmo sem saber qual inversa generalizada o colega usou.
\end{enumerate}
\end{questao}

\end{document}
